% sections/01_problem_framing.tex  —  Problem Framing
% \input{} into main.tex; do NOT wrap in \begin{document}.

\section{Problem Framing}
\label{sec:problem}

\subsection{The $S_{11}$ Reflection Coefficient}
\label{sec:problem:s11}

In microwave network analysis the scattering parameter $S_{11}$
characterises the ratio of the complex voltage wave reflected from a
one-port device to the complex voltage wave incident upon
it~\cite{pozar2011microwave}:
\begin{equation}
    S_{11}(f) \;=\; \frac{V^{-}(f)}{V^{+}(f)},
    \label{eq:s11_def}
\end{equation}
where $V^{+}$ is the incident travelling wave and $V^{-}$ is the
reflected travelling wave at frequency~$f$.  Because $S_{11}$ is a
complex quantity, it can be represented in either polar form (magnitude and
phase) or Cartesian form (real and imaginary components). While human engineers
typically prefer the polar form to visualise return loss in \si{\deci\bel},
the raw data files in this project store $S_{11}$ in precisely this
polar format ($|S_{11}|$ in~dB, $\angle S_{11}$ in degrees).  To
construct a representation that is suitable for gradient-based learning,
the pipeline converts each trace into a \textbf{three-channel hybrid
tensor} $[\,|S_{11}|_{\text{dB}},\;\text{Re}[S_{11}],\;
\text{Im}[S_{11}]\,]$, where the Cartesian components are derived from
the native polar values.  This approach preserves the high dynamic range
of the dB magnitude while providing mathematically continuous channels
that avoid the ``phase wrapping'' problem---the artificial discontinuity
that occurs when phase angles wrap from $\SI{+180}{\degree}$ to
$\SI{-180}{\degree}$.

\subsection{Physical Origin of Class Separability}
\label{sec:problem:physics}

The sensor under study is a single-port microwave structure whose
electromagnetic near-field interacts with the material present on its
surface.  When the surface condition changes—from a dry dielectric
substrate to a thin water film or an ice layer—the effective complex
permittivity $\varepsilon^{*} = \varepsilon' - j\varepsilon''$ of the
medium shifts.  This perturbation modifies the resonant frequencies, the
quality factors, and the coupling coefficients of the sensor's radiating
or guided-wave structure, producing measurable changes in the
$S_{11}(f)$ trace across the \SIrange{8}{12}{\giga\hertz} band.

Specifically:
\begin{itemize}
    \item \textbf{Dry surface} (\texttt{dry}):  The sensor is loaded
        only by the substrate and ambient air.  The $S_{11}$ trace
        exhibits the sensor's baseline resonance profile.
    \item \textbf{Wet surface} (\texttt{wet}):  A liquid water film
        introduces a material with high real permittivity
        ($\varepsilon'_{\text{water}} \approx 50$--$80$ at X-band
        frequencies) and significant dielectric loss.  This shifts
        resonances downward in frequency, broadens them, and deepens or
        shallows dips depending on coupling geometry.
    \item \textbf{Ice-covered surface} (\texttt{ice}):  Ice has a
        substantially lower permittivity ($\varepsilon'_{\text{ice}}
        \approx 3.2$) and much lower loss than liquid water.  The
        resulting $S_{11}$ perturbation is smaller in magnitude than the
        wet case but distinct from the dry baseline.
\end{itemize}

The classification task exploits these physically grounded spectral
differences.  However, the degree of separability is an empirical
question that depends on the sensor geometry, the thickness and
uniformity of the surface layer, and the measurement noise floor—none
of which are known \emph{a priori} from the data files alone.

\subsection{Classification Goal}
\label{sec:problem:goal}

The formal objective is to learn a mapping
\begin{equation}
    h \;:\; \mathbf{x}(f) \;\longmapsto\; y \;\in\; \{\texttt{dry},\;
        \texttt{wet},\;\texttt{ice}\},
    \label{eq:classifier}
\end{equation}
where $\mathbf{x}(f)$ is the feature representation derived from a
single $S_{11}$ frequency sweep and $y$ is the categorical surface-state
label.  This is a \emph{closed-set, single-label, three-class}
classification problem.  Each trace is assumed to originate from exactly
one surface condition, and no ``unknown'' or out-of-distribution class is
modelled at this stage.

The mapping~$h$ must generalise from the \num{1603} available labelled
traces to previously unseen measurements acquired under nominally
identical conditions.  The small sample size is a dominant methodological
constraint: it limits the complexity of models that can be reliably
validated and increases the risk of overfitting, particularly for
high-dimensional deep architectures.

\subsection{Edge Deployment Context}
\label{sec:problem:edge}

The trained classifier is intended for deployment on
resource-constrained edge devices positioned at or near the sensor
installation site.  While the specific target hardware has not been
finalised, typical edge platforms impose the following constraints:

\begin{itemize}
    \item \textbf{Memory budget.}  Model weights and inference buffers
        must fit within tens of kilobytes to a few megabytes of RAM,
        depending on the platform.
    \item \textbf{Compute budget.}  Inference must complete within
        milliseconds to low hundreds of milliseconds.  Floating-point
        units may be absent or limited to single precision; integer-only
        or fixed-point arithmetic may be required.
    \item \textbf{Power budget.}  Continuous or near-continuous inference
        must be sustainable within the power envelope of the deployed
        system.
    \item \textbf{Software ecosystem.}  Runtime environments are
        typically limited to lightweight inference engines such as
        TensorFlow Lite~\cite{tensorflow_lite} or ONNX
        Runtime~\cite{onnx_runtime}, rather than full training
        frameworks.
\end{itemize}

These constraints mean that a model achieving the highest validation
accuracy in a desktop experiment is not automatically the best candidate
for deployment.  Model size, quantisability, and inference throughput are
co-equal selection criteria alongside predictive performance.  The
methodology in this report is designed to produce comparable metrics
across all three dimensions.

\subsection{Scope and Boundaries}
\label{sec:problem:scope}

To maintain focus and methodological rigour, this report adheres to the
following explicit scope boundaries:

\begin{enumerate}
    \item \textbf{Single-port only.}  The analysis is restricted to
        $S_{11}$ measurements from a one-port network.  Multi-port
        extensions involving $S_{21}$, $S_{12}$, or $S_{22}$ are out of
        scope.
    \item \textbf{No time-series or drift modelling.}  Each trace is
        treated as an independent, identically distributed sample.
        Temporal ordering of acquisitions, sensor drift, and sequential
        decision-making are not considered.
    \item \textbf{No experimental results.}  This document defines the
        methodological framework—representations, preprocessing
        pipelines, candidate algorithms, and evaluation protocols.
        Trained model performance numbers are deferred entirely to the
        companion experiment report.
    \item \textbf{No manuscript-ready claims.}  Because no experiments
        have been executed, no claims of superiority, state-of-the-art
        performance, or benchmark comparisons are made.
    \item \textbf{Closed-set assumption.}  The classifier is designed for
        the three known classes only.  Detection of novel or anomalous
        surface states (out-of-distribution detection) is noted as a
        future direction but is not addressed here.
\end{enumerate}
